
\documentclass[10pt, conference, compsocconf]{IEEEtran}

\usepackage[utf8]{inputenc}
\usepackage[T1]{fontenc}
\usepackage[brazilian]{babel}
\usepackage{graphicx,url}
\usepackage{placeins}

\usepackage{amsmath}
%\usepackage{cite}
\usepackage{xcolor}
%\Elaborado por Fabrício Tonetto Londero
\usepackage{listings}

\usepackage[fixlanguage]{babelbib}
\selectbiblanguage{brazil}

%\usepackage[backend=biber]{biblatex}
\usepackage[ natbib=true, style=numeric,sorting=none]{biblatex} %ordem de ocorrencia da referencia

\addbibresource{references.bib}

\hyphenation{op-tical net-works semi-conduc-tor}

\begin{document}

\title{Título do trabalho}

\author{\IEEEauthorblockN{Aluno, Orientador}
\IEEEauthorblockA{Curso de [Nome do Curso]\\
UFN - Universidade Franciscana\\
Santa Maria - RS\\
aluno@ufn.edu.br, 
orientador@ufn.edu.br}
}

\maketitle

\begin{abstract}
Lorem ipsum dolor sit amet, consectetur adipiscing elit. Proin vitae erat velit. Donec nec euismod velit. In ut ipsum lorem. Nulla vitae metus ut nibh mollis sodales in eu nulla. Pellentesque enim metus, interdum eget iaculis vitae, posuere non massa. Pellentesque sit amet libero vitae dolor dapibus porta. Morbi id eros eu diam posuere fermentum. \textit{Phasellus ac orci} vel nisl elementum vulputate. Quisque vitae malesuada neque. Duis posuere mi nec purus tincidunt sollicitudin. Ut cursus luctus condimentum. Aliquam ut tortor at nulla mollis efficitur.

\end{abstract}


\textbf{\textit{Palavras-chave} : palavra 1; palavra 2; palavra 3; }


\IEEEpeerreviewmaketitle

\section{Introdução
Introduzir.
Os jogos sao tops demais. E temos que usar eles para algo alem de jogar e resenhar. Que tal aprender?
\subsection{Objetivos}

Propor o design de um jogo capaz de suprir as demandas apresentadas neste texto. Baseado em SDT (Self-deternmination Theory), Aprendizagem tangencial e as teorias de Ausubel, quatno ao Scaffolding (o conhecimento precisa ter uma base para ser construido).

\subsection{Objetivos secundarios}
Bah.

\subsection{Estrutura do trabalho}
Vai ter o que estiver no titulo. E vou expandir na parte das ferramentas. E da metodologia, e do design, ja que vou focar bastante nessa parte

\FloatBarrier

\begin{figure}[!h]
\centering
\includegraphics[width=3.3in]{images/Overwatch.png}
\caption{Legenda da Figura \cite{FreyaHolmer}.} 
\label{figuraOverwatch}
\end{figure}

\FloatBarrier

\section{Secao 1 ou Referencial teorico}
Falar sobre: o que eh, validar. Como vai ser usado, validar.

\subsection{SDT}
Vai ser o mais importante?
Tem muita pesquisa que ja valida. Ja foi aplicado em multiplas pesquisas, os manos que criaram seguem acreditando e escrevendo sobre. Vai ser usado para guiar o desenho do design (?) das mecanicas do jogo. Trazer algum projeto que tenha seguido esse principio. 


\subsection{Scafolding do Ausubel ou Trecos de sustentacao} 
Apesar de usar o nome do cara, eu nao vou usar a teoria dele. Vou usar uma versao 100% atualizada da ideia dele, em que, o conhecimento usado de scafodling, para aprender o aprendizado, nao eh jogado fora. Acho que eh recente o trabalho dos manos, entao tem pouca coisa. 
Vou usar para dar ao jogador algo nao-abstrato para ancorar o conhecimento. O atomo de hidrogenio vai ser um peixe/tijolo por ser a base para todos e ser fraquinho. Bah. Nao sei se ja foi aplicado em jogo.

\subsection{Aprendizgem tangencial}
Apesar de ser recente, tem bastante estudo sobre. Nao vai ser dificil, por ser algo bem exclusivo de jogo.
Vai ser aplicado como foi aplicado em todos os outros, provendo entradas numa Pokedex, com dados sobre os elementos ou as interacoes. Deve ter jogo que ja validou isso.  

\subsect{Coisas de jogo}
Elaborar sobre o Game design, estilos graficos.

\section{Secao 2 Ferramentas ou Open-source e gratuito}

Listar o que vou usar, como vou usar e porque usar. Acho que a filosofia do Game Design cai aqui dentro.

\subsection{Engine}
Usar o Godot. Gosto, de graca e open source. As ferramentas que ele possui sao boas. Menos documentacao sobre.

\subsection{Treco de graficos}
Blender, Krita, Gimp, Inkscape.

\subsection{IDE}
Neo Vim. Atalhos.


\subsection{Trabalhos correlatos}
Jogo de quimica, pokemao/digimon.
Algum jogo com Aprendizgem Tangencial, um que me tenha feito sair do jogo pra procurar sobre.
Jogos pra jogar: Digimon World, Pokemon Rogue
\section{Conclusão}

Vai dar certo, no jeitao dos rapazes
